\color{unidarkgreen}{\subsection{Контроль изоляции стороны НН}}
\color{black}

\begin{enumerate}[label=\arabic{section}.\arabic{subsection}.\arabic*,  align=left, labelsep=4pt, leftmargin=0pt, itemindent=57pt]

\item
Для информирования персонала при возникновении однофазного замыкания на землю
предусмотрен сигнальный орган контроля изоляции, контролирующий значение утроенного напряжения
нулевой последовательности (3U0). 

\item
Значение срабатывания по напряжению нулевой последовательности, определяемое параметром
\textbf{«3U0ср КИ НН»}. Для исключения ложного формирования сигнализации при возникновении
неисправности цепей напряжения предусмотрен блокирующий орган, включенный на напряжение
обратной последовательности. Значение блокировки по напряжению обратной последовательности
задается параметром \textbf{«U2бл КИ НН»}. Действие на сигнализацию формируется с задержкой по времени
\textbf{«Тср КИ НН»}. 

\begin{figure}[H]
\centering
\includegraphics[width=0.8\textwidth,height=0.8\textheight,keepaspectratio]{\fbpath/053.0961 КИ НН/img/KINN.pdf}
\captionof{figure}{Функциональный блок КИ НН}\label{pic:KINN}
\end{figure}


\small
\begin{longtable}{|>{\centering\arraybackslash}m{5.3cm}|>{\centering\arraybackslash}m{4.2cm}|>{\centering\arraybackslash}m{1cm}|>{\centering\arraybackslash}m{3.3cm}|>{\centering\arraybackslash}m{1.8cm}|}
% Первая страница: полный заголовок
\caption{Перечень уставочных параметров функции КИ НН\hfill\vspace{-0.5\baselineskip}}\label{desc:tbl1}\\
\hline
\rowcolor{gray!30}
Наименование уставки & Диапазон & Шаг & Значение по умолчанию & Единица измерения \\
\hline
\endfirsthead

% Промежуточные страницы: «Продолжение таблицы»
\caption*{\hspace{3pt}\emph{Продолжение таблицы \ref{desc:tbl1}}\hfill\ }\vspace{-0.5\baselineskip} \\
\hline
\rowcolor{gray!30}
Наименование уставки & Диапазон & Шаг & Значение по умолчанию & Единица измерения \\
\hline
\endhead

%===m>ДЗАТ|-
%===m
\end{longtable}
\normalsize	

\end{enumerate}